\scanpage{68}
\mspage{67}

on trouve
\[
\Sigma_{\unsure{L|\mathcal{L}}}=\unsure{\mathscr{P}|\mathcal{P}}(L),
\]
l'opération $C_L$ étant simplement le passage au complémentaire. Donc dans ce cas, le support d'une partie $A$ de $\unsure{\mathcal{F}|\mathscr{F}}$ (ou de $M$, quand on part des \unsure{Anatos|Aratos}) ; identifié à une partie de $L$, n'est autre que $\operatorname{omb}_L(A)=\operatorname{omb}(A)$. Donc dans ce cas, la définition revient simplement à poser
\[
\begin{aligned}
\operatorname{supp}A&=\operatorname{omb}(A)\overset{\mathrm{def}}{=}\bigcup_{F\in A}\operatorname{omb}(F)\\
\operatorname{cosupp}A&=\unsure{\complement|C}_L\operatorname{supp}A
\qquad\begin{gathered}\text{(complémentaire}\\\text{ordinaire)}\end{gathered}
\end{aligned}
\]

Ceci tient au fait que la relation cruciale de disjonction devient simplement
\diagram{\begin{tikzcd}[ampersand replacement=\&, column sep=small, row sep=large]
F\parallel G\quad\Longleftrightarrow\quad\operatorname{omb}F\cap\operatorname{omb}G={} \& \varnothing \arrow[d, bend right=15] \& {,} \\
\text{relation de disjonction} \& \text{ensembliste} \& \text{sur les}
\end{tikzcd}}
omb (\unsure{rebaptisés} supports pour la circonstance). Dans ce cas, les propriétés p.63,64 se \unsure{résolvent}, on a
\[
\left\{\begin{aligned}
A\subset B&\quad\Longleftrightarrow\quad A\cap\underset{\substack{\uparrow\\\text{``compl.''}}}{\unsure{\mathbb{C}|\complement}}(B)=\varnothing_\Sigma\\
A\cap\bigvee_i B_i&=\bigvee_i A\cap B_i\qquad\text{dans }\Sigma\quad-
\end{aligned}\right.
\]
